\section*{Chapter \ref{ch:mx:why-there-is-a-spread} --- Why There Is a Spread}

\begin{solution}{exo:mx:why-there-is-a-spread:1}
$s=2\sqrt{0.0004}=0.04$: four cents.
\end{solution}

\begin{solution}{exo:mx:why-there-is-a-spread:2}
With $\delta=\tfrac12$: $a=99+2\times\tfrac{1.2}{2}=100.2$ and $b=99+2\times\tfrac{0.8}{2}=99.8$; the spread is $\mu\times2=0.4$.
\end{solution}

\begin{solution}{exo:mx:why-there-is-a-spread:3}
$\lambda=2/(2\times8)=0.125$, $\beta=8/2=4$, profit $\tfrac12\times2\times8=8$, posterior variance $4/2=2$.
\end{solution}

\begin{solution}{exo:mx:why-there-is-a-spread:4}
$\gamma\sigma^2T=0.1$: $a=50+0.1\,(0.5+3)=50.35$ and $b=50-0.1\,(0.5-3)=50.25$. Being short, she wants to buy: her reservation price is
$50+0.3=50.3$, and both quotes sit around it, above the value.
\end{solution}

\begin{solution}{exo:mx:why-there-is-a-spread:5}
The belief becomes the old ask on the unit scale: $\delta'=0.6$. New quotes: $b=100.00$ and $a=100.38$; the spread narrows from 0.40 to 0.38 and both quotes
move up.
\end{solution}

\begin{solution}{exo:mx:why-there-is-a-spread:6}
The insider scales his trading with the noise: $\beta=\sigma_u/\sqrt{\Sigma_0}$ doubles, the signal-to-noise ratio of the order flow is unchanged, and so is
the share of information in the price (one half). What changes is $\lambda$ (halved: a deeper market) and the insider's profit (doubled).
\end{solution}

\begin{solution}{exo:mx:why-there-is-a-spread:7}
On one sample (seed 41) PIN is 0.149 with the swings, 0.139 after normalising, against 0.098 on the same sample without swings (true 0.107). Normalising removes
the common factor but also the extra volume of informed days, which the model uses to spot them, and the rescaled counts are no longer Poisson (their
variance is not their mean), so the likelihood is misspecified in another way.
\end{solution}

\begin{solution}{exo:mx:why-there-is-a-spread:8}
PIN is a model-based share of trades, not a measure of informed losses: common swings in activity, not information, can raise it (the chapter's simulation
adds 0.15 for swings of standard deviation 0.5), and it says nothing about how much each informed trade costs the market maker. Compare adverse-selection
costs directly (chapter~5) before concluding.
\end{solution}

\begin{solution}{pb:mx:why-there-is-a-spread:1}
\textbf{1.} A constant half-spread, independent trade signs and a random-walk value; 0.88 ticks on the simulated hour.
\textbf{2.} The simulated trade signs are positively autocorrelated (0.29 at lag one): the bounce is partly cancelled, and the estimator falls short of
the quoted 1.08 ticks.
\textbf{3.} $a=v+\gamma\sigma^2T(Q/2-q)$, $b=v-\gamma\sigma^2T(Q/2+q)$; spread $\gamma\sigma^2TQ$.
\textbf{4.} Both move down by $\gamma\sigma^2Tq$: the dealer wants to sell and to avoid buying more.
\textbf{5.} $\Pr(v_H\mid\text{buy})=\tfrac{(1+\mu)/2}{(1+\mu)/2+(1-\mu)/2}=\tfrac{1+\mu}{2}$, so $a=v_L+(v_H-v_L)\tfrac{1+\mu}{2}$.
\textbf{6.} 0.3.
\textbf{7.} 148, 19 and 6 trades.
\textbf{8.} One half and three quarters.
\textbf{9.} See the proof of the proposition: $\beta=1/(2\lambda)$ from the insider's problem, $\lambda=\beta\Sigma_0/(\beta^2\Sigma_0+\sigma_u^2)$ from the
market maker's, hence $\beta=\sigma_u/\sqrt{\Sigma_0}$, $\lambda=\sqrt{\Sigma_0}/(2\sigma_u)$.
\textbf{10.} Half: $\mathrm{Var}(v\mid p)=\Sigma_0/2$.
\textbf{11.} 0.250, 0.498 and 0.991 (against 0.25, 0.5 and 1).
\textbf{12.} The \omterm{def:mx:why-there-is-a-spread:informed}{noise traders}, whose orders execute at prices moved by the insider's.
\textbf{13.} $\alpha\mu/(\alpha\mu+\varepsilon_b+\varepsilon_s)=24/224=0.107$.
\textbf{14.} 0.006.
\textbf{15.} 0.085, 0.153 and 0.207.
\textbf{16.} \emph{Named result}: with a processing cost of 0.05 per trade and a value range of 1, adverse selection is $\mu/(\mu+0.1)$ of the spread
(one half at $\mu=0.1$, three quarters at $\mu=0.3$); PIN estimated on days with no information rises from 0.006 to 0.085, 0.153 and 0.207 as common activity
swings grow to standard deviations of 0.25, 0.5 and 0.75, and an informed market's 0.107 is read as 0.195 with the largest.
\textbf{17.} A day with more buys and sells than usual fits the event-day components better than the no-event component, whose rates are fixed; the model has
no other way to explain dispersion in daily totals.
\textbf{18.} That the component of PIN related to illiquidity, not the one related to asymmetric information, is priced in the cross-section of returns.
\textbf{19.} Model the common factor explicitly (time-varying uninformed rates, symmetric order-flow shocks), and measure information by what follows
trades: price impact that persists (chapter~5's permanent component) rather than volume.
\textbf{20.} In a liquid stock with a one-tick spread, mostly adverse selection and processing costs net of rebates; inventory costs are small because
positions are turned over in seconds.
\end{solution}

\begin{solution}{iq:mx:why-there-is-a-spread:1}
Order processing (evidence: spreads fall with fees and technology costs; Roll's bounce), inventory (quotes move against dealer positions; spreads widen
with volatility and trade size), adverse selection (prices move in the direction of trades and stay there; spreads widen before announcements).
\iqlookfor{the three names and a mechanism-level piece of evidence for each.}
\end{solution}

\begin{solution}{iq:mx:why-there-is-a-spread:2}
$\Pr(\text{buy}\mid1)=(1+\mu)/2$, $\Pr(\text{buy}\mid0)=(1-\mu)/2$, so $\Pr(1\mid\text{buy})=(1+\mu)/2$ and, with zero profit, $a=(1+\mu)/2$; the bid is
$(1-\mu)/2$ and the spread $\mu$. \iqlookfor{Bayes' rule and the zero-profit condition, stated before computing.}
\end{solution}

\begin{solution}{iq:mx:why-there-is-a-spread:3}
$\lambda$ halves (impact falls with the depth that noise provides), $\beta$ doubles, the insider's profit doubles, and the price reveals the same half of his
information. \iqlookfor{that informativeness does not depend on noise, because the insider scales up.}
\end{solution}

\begin{solution}{iq:mx:why-there-is-a-spread:4}
Lower both quotes (skew): a higher chance of selling and a lower chance of buying more, and possibly widen if the position is near its limit; hedge in a
correlated instrument if one exists. \iqlookfor{the reservation-price idea and the trade-off with adverse selection when skewing.}
\end{solution}

\begin{solution}{iq:mx:why-there-is-a-spread:5}
The first-order autocovariance of price changes was positive, not negative: trade signs are autocorrelated (order splitting), prices trend over the
sampling interval, or the data mixes quotes from several venues. Use quote data if available, sample at a coarser interval, or use an estimator that does
not rely on the bounce (Corwin--Schultz, Abdi--Ranaldo; chapter~5). \iqlookfor{knowing why the estimator fails, not only that it does.}
\end{solution}

\begin{solution}{iq:mx:why-there-is-a-spread:6}
Decompose it: measure the effective spread and the realised spread at a horizon where price moves have settled; their difference is the price impact that
stays, the adverse-selection part; or fit a structural decomposition (Huang--Stoll, MRR) or a VAR. \iqlookfor{the effective--realised decomposition and
the choice of horizon, both chapter~5.}
\end{solution}
